\begin{enumerate}
	\item Au bout de 24 mois, le prix payé par Lola est :
	\begin{itemize}
		\item \textbf{Offre A} : \quad $175+16\times 24=559$;
		\item \textbf{Offre B} : \quad $23\times 24=552$
	\end{itemize}

	$552 < 559$, \quad donc l'offre B est plus intéressante au bout de 24 mois.

	\item \begin{enumerate}
		\item Comme $x$ représente le nombre de mois, on doit prendre le tarif par mois et le multiplier par $x$, et puis, on doit ajouter le prix payé à l'achat, éventuellement.

		La fonction $f$ correspond à l'offre \textbf{A} et la fonction $g$ à l'offre \textbf{B}.

		\item Puisque $f(x)$ est le prix payé au bout de $x$ mois avec l'offre A et $g(x)$ celui payé avec l'offre B, pour répondre à la question, on cherche à résoudre l'équation :
		\begin{align*}
		f(x)&=g(x)\\
		175+16x&=23x\\
		16x-23x&=-175\\
		-7x&=-175\\
		x&=\dfrac{-175}{-7}\\
		x&=25
		\end{align*}

		La solution de l'équation est 25.

		Cela signifie qu'au bout de 25 mois, les deux formules donnent le même prix.

		\item $25>24$,\quad donc on n'est plus dans la période d'engagement.
	\end{enumerate}
\end{enumerate}
